

\chapter{Technical Analysis of MATLAB Function: \texttt{linksection.m}}
\author{}
\date{}


\maketitle

\section{Overview}

Based on the provided source code, \texttt{linksection.m} currently serves as a functional placeholder for a geometric transformation routine.

\begin{itemize}
    \item \textbf{Main Purpose}: The file is intended to implement a "link section" operation that transforms a given complex-link structure \texttt{(g, phi)} into a new one \texttt{(g\_, phi\_)}.
    \item \textbf{Type of Analysis}: The function is designed to facilitate morphological operations on complex-links, as described in specific topological literature.
    \item \textbf{Mathematical/Theoretical Context}: The theoretical basis for this file is situated in \textbf{Discrete Geometry} and \textbf{Complex Analysis}, specifically focusing on the winding numbers of complex-links and their applications.
\end{itemize}

\section{Step-by-Step Code Analysis}

\subsection{Function Definition and Interface}
\textbf{Explanation}: This section establishes the function signature. It defines 
\texttt{linksection} as a routine that accepts a complex-link structure defined by a set of complex points 
\texttt{g} and an endomorphism (endomap) 
\texttt{phi}. It is designed to return a modified set of points 
\texttt{g\_} and a new connectivity map 
\texttt{phi\_}.

\begin{lstlisting}
[g_,phi_]=linksection(g,phi)
\end{lstlisting}

\subsection{Implementation Status and Reference}
\textbf{Explanation}: The code includes a timestamp and a note stating that the internal logic is not yet implemented. It provides a direct academic reference to the paper "On the winding number of a complex-link and its applications to morphological operations" by N. Martins-Ferreira (2022), citing Section 2.4 on page 23 as the basis for the intended logic.

\begin{lstlisting}
% 27May2025 -- not implemented yet
%
% Given a link structure (g,phi) returns a new link structure (g_,phi_)  
% as described in section 2.4 of page 23 of the paper: N. Martins-Ferreira, 
% "On the winding number of a complex-link and its applications to 
% morphological operations", Scripta-Ingenia 11 (December) (2022) 20--26.;
\end{lstlisting}

\section{Expected Outputs and Visuals}



As the logic is currently not implemented, the function does not produce numerical or graphical outputs. However, based on the referenced academic paper:

\begin{itemize}
    \item \textbf{Complex-Link Visualization}: A plot of a complex-link \texttt{(g, phi)} would display discrete complex points in a 2D plane connected by directed edges.
    \item \textbf{Morphological Result}: The "sectional" operation is used to analyze or transform the topological properties of these shapes.
\end{itemize}

The final output of an implemented version would show a transformed set of vertices on the \textbf{Real} and \textbf{Imaginary} axes, illustrating the morphological refinement of the input shape.

